% ═══════════════════════════════════════════════════════════════════════════
%  فصل ۹ — استعمار، پسااستعمار و تجزیه‌ی مرزها
% ═══════════════════════════════════════════════════════════════════════════
\chapter{استعمار، پسااستعمار و تجزیه‌ی مرزها}
\label{ch:colonial}

\begin{tcolorbox}[mybox, title={چکیده‌ی فصل}]
این فصل به بررسی رابطه‌ی نجات‌طلبی با فرایندهای استعماری و پسااستعماری می‌پردازد. سه محور اصلی: (۱) استعمار اولیه و ثانویه و نقش «دعوت» یا «رضایت» بومیان؛ (۲) استعمارزدایی و نجات‌طلبی معکوس (جنبش‌های آزادی‌بخش)؛ (۳) تجزیه‌ی مرزها و تشکیل کشورهای جدید — کجا موفق و کجا ناموفق. دسته‌بندی مفهومی فعالیت‌های میسیونری، پان‌ایدئولوژیک و امپریالیستی ارائه می‌شود.
\end{tcolorbox}

% ─────────────────────────────────────────────────────────────────────────
\section{استعمار و مسئله‌ی «رضایت»: آیا بومیان دعوت کردند؟}
\label{sec:ch09-colonial-consent}
% ─────────────────────────────────────────────────────────────────────────

\subsection{استعمار اولیه: «کشف»، فتح و تسهیل‌گری محلی}
\label{subsec:ch09-early-colonialism}

فرایند استعمار اروپایی (سده‌های ۱۵–۱۹ میلادی) به‌ندرت با «دعوت» صریح بومیان آغاز شد. با این حال، \concept{تسهیل‌گری محلی} نقش مهمی داشت: در بسیاری از موارد، بخشی از نخبگان بومی — به‌دلیل رقابت‌های داخلی — با استعمارگران همکاری کردند.%
\footnote{\lr{Robinson, Ronald. "Non-European Foundations of European Imperialism." In \textit{Studies in the Theory of Imperialism}, ed. Owen and Sutcliffe. Longman, 1972, pp.~117--140.}}

\thinker{رونالد رابینسون} (\lr{Ronald Robinson}) نظریه‌ی «بنیادهای غیراروپایی امپریالیسم» را مطرح کرده و استدلال می‌کند که استعمار بدون \concept{همکاری نخبگان محلی} (\lr{collaborative mechanism}) ممکن نبود:

\begin{longtable}{
  >{\raggedleft\arraybackslash\bfseries}p{2.5cm}
  >{\raggedleft\arraybackslash}p{3.5cm}
  >{\raggedleft\arraybackslash}p{3.5cm}
  >{\raggedleft\arraybackslash}p{3.5cm}}
\caption{نمونه‌های تسهیل‌گری محلی در استعمار}\label{tab:ch09-collaboration}\\
\toprule
\rowcolor{PrimaryDeep}
\textcolor{white}{\textbf{منطقه}} &
\textcolor{white}{\textbf{تسهیل‌گران محلی}} &
\textcolor{white}{\textbf{انگیزه}} &
\textcolor{white}{\textbf{نتیجه}} \\
\midrule
\endfirsthead
\toprule
\rowcolor{PrimaryDeep}
\textcolor{white}{\textbf{منطقه}} &
\textcolor{white}{\textbf{تسهیل‌گران محلی}} &
\textcolor{white}{\textbf{انگیزه}} &
\textcolor{white}{\textbf{نتیجه}} \\
\midrule
\endhead
\bottomrule
\endlastfoot

هند &
راجاها و نوّاب‌های محلی &
رقابت با رقبای داخلی &
سلطه‌ی شرکت هند شرقی%
\footnote{\lr{Bayly, C.A. \textit{Indian Society and the Making of the British Empire}. Cambridge UP, 1988.}} \\
\rowcolor{NeutralLight}
آفریقای غربی &
رؤسای قبایل ساحلی &
تجارت و سلاح &
استعمار کامل \\
آمریکای مرکزی &
تلاکسکالتک‌ها (متحد کورتز) &
دشمنی با آزتک‌ها &
نابودی آزتک‌ها، استعمار اسپانیا%
\footnote{\lr{Restall, Matthew. \textit{Seven Myths of the Spanish Conquest}. Oxford UP, 2003.}} \\
\rowcolor{NeutralLight}
اندونزی &
سلاطین محلی جاوه &
تحکیم قدرت محلی &
سلطه‌ی هلند \\
آفریقای شرقی &
سلطان زنگبار &
تجارت و حمایت &
حمایت‌خواهی بریتانیا \\

\end{longtable}

\subsection{فعالیت‌های میسیونری: تسهیل‌گری فرهنگی استعمار}
\label{subsec:ch09-missionary}

\concept{فعالیت‌های میسیونری} (اعم از مسیحی، و بعدها ایدئولوژیک) نقش مهمی در تسهیل پذیرش استعمار داشتند. \thinker{جان و جین کاماروف} (\lr{John and Jean Comaroff}) در پژوهش تأثیرگذار خود نشان داده‌اند که میسیونرها «ذهن» بومیان را پیش از ورود ارتش فتح می‌کردند:

\begin{histQuote}{کاماروف و کاماروف}
«میسیونرها پیش‌قراولان استعمار بودند — نه به‌معنای توطئه‌ی آگاهانه، بلکه به این معنا که با تغییر ارزش‌ها، عادات و ساختار اجتماعی بومیان، مقاومت فرهنگی در برابر استعمار را تضعیف می‌کردند.»%
\footnote{\lr{Comaroff, John L., and Jean Comaroff. \textit{Of Revelation and Revolution: Christianity, Colonialism, and Consciousness in South Africa}. U of Chicago P, 1991, vol.~1, p.~4.}}
\end{histQuote}

% ─────────────────────────────────────────────────────────────────────────
\section{طبقه‌بندی: میسیونری، پان‌ایدئولوژیک یا امپریالیستی؟}
\label{sec:ch09-classification}
% ─────────────────────────────────────────────────────────────────────────

\begin{tcolorbox}[defbox, title={سه‌گانه‌ی مفهومی: تسهیل‌گری بیگانه}]
فعالیت‌هایی که مداخله‌ی بیگانه را «نرم و پذیرفتنی» می‌کنند، در سه دسته قابل تقسیم‌اند:

\begin{enumerate}[nosep, label=\textbf{\arabic*.}]
  \item \textbf{میسیونری-مذهبی:} تبلیغ دین (مسیحیت، اسلام) که پیوند هم‌کیشی فرامرزی ایجاد کرده و مرز «خودی/بیگانه» را بازتعریف می‌کند.
  \begin{itemize}[nosep]
    \item مسیحی‌سازی آفریقا و آسیا (سده‌های ۱۶–۲۰)
    \item صدور انقلاب اسلامی ایران (۱۳۵۷–کنون)
  \end{itemize}

  \item \textbf{پان‌ایدئولوژیک:} ایدئولوژی‌های فراملی که هویت مشترک (نژادی، زبانی، طبقاتی) را بر هویت ملی ترجیح می‌دهند:
  \begin{itemize}[nosep]
    \item \concept{پان‌اسلاویسم:} روسیه به‌عنوان «حامی اسلاوها» — ابزار نفوذ در بالکان%
    \footnote{\lr{Kohn, Hans. \textit{Pan-Slavism: Its History and Ideology}. 2nd ed. Vintage, 1960.}}
    \item \concept{پان‌ترکیسم:} ترکیه/عثمانی به‌عنوان «رهبر ترکان» — ابزار نفوذ در قفقاز و آسیای مرکزی%
    \footnote{\lr{Landau, Jacob M. \textit{Pan-Turkism: From Irredentism to Cooperation}. 2nd ed. Indiana UP, 1995.}}
    \item \concept{پان‌عربیسم:} ناصر/بعث به‌عنوان «متحدکننده‌ی عرب‌ها» — ابزار مداخله در کشورهای عربی%
    \footnote{\lr{Dawisha, Adeed. \textit{Arab Nationalism in the Twentieth Century}. Princeton UP, 2003.}}
    \item \concept{انترناسیونالیسم کمونیستی:} شوروی به‌عنوان «رهبر پرولتاریای جهان» — ابزار نفوذ در سراسر دنیا
  \end{itemize}

  \item \textbf{امپریالیستی صرف:} مداخله بدون هیچ پیوند هم‌کیشی، هم‌نژادی یا هم‌ایدئولوژیکی — صرفاً بر اساس برتری نظامی و منفعت اقتصادی:
  \begin{itemize}[nosep]
    \item استعمار بریتانیا در هند
    \item استعمار بلژیک در کنگو%
    \footnote{\lr{Hochschild, Adam. \textit{King Leopold's Ghost}. Houghton Mifflin, 1998.}}
    \item حمله‌ی مغول به ایران
  \end{itemize}
\end{enumerate}
\end{tcolorbox}

\begin{figure}[htbp]
\centering
\begin{tikzpicture}[
  every node/.style={font=\small},
  circle1/.style={circle, draw=AccentTeal, fill=BgTeal!50, minimum size=4cm, opacity=0.6},
  circle2/.style={circle, draw=AccentAmber, fill=BgWarm!50, minimum size=4cm, opacity=0.6},
  circle3/.style={circle, draw=AccentRed, fill=BgFail!50, minimum size=4cm, opacity=0.6},
]

% سه دایره‌ی ون
\node[circle1] (miss) at (0,0) {};
\node[circle2] (pan) at (3,0) {};
\node[circle3] (imp) at (1.5,-2.5) {};

% برچسب‌ها
\node[font=\small\bfseries, color=AccentTeal] at (-1.2,0.8) {\rl{میسیونری}};
\node[font=\small\bfseries, color=AccentAmber] at (4.2,0.8) {\rl{پان‌ایدئولوژیک}};
\node[font=\small\bfseries, color=AccentRed] at (1.5,-4) {\rl{امپریالیستی}};

% تقاطع‌ها
\node[font=\scriptsize, text width=1.5cm, align=center] at (1.5,0.3)
  {\rl{صدور انقلاب}};
\node[font=\scriptsize, text width=1.5cm, align=center] at (0.3,-1.4)
  {\rl{استعمار مسیحی}};
\node[font=\scriptsize, text width=1.5cm, align=center] at (2.7,-1.4)
  {\rl{پان‌اسلاویسم + توسعه‌طلبی}};
\node[font=\scriptsize, text width=1.5cm, align=center] at (1.5,-1.2)
  {\rl{ترکیب هر سه}};

\end{tikzpicture}
\caption{نمودار ون: سه‌گانه‌ی تسهیل‌گری بیگانه}
\label{fig:ch09-venn}
\end{figure}

\begin{tcolorbox}[defbox, title={پاسخ به پرسش کلیدی: آیا نجات‌طلبی فقط میان ملت‌های «آشنا» ممکن است؟}]
بر اساس شواهد تاریخی:

\begin{itemize}[nosep, rightmargin=1em]

  \item \textbf{نجات‌طلبی واقعی} (دعوت آگاهانه از بیگانه) عمدتاً در میان ملت‌هایی رخ می‌دهد که \concept{پیوند ارگانیک} (مذهبی، زبانی، نژادی، ایدئولوژیک) با مداخله‌گر دارند: گرجی‌ها و روسیه (هم‌کیشی)، لهستانی‌ها و ناپلئون (ایدئولوژی)، شیعیان عراق و ایران (هم‌مذهبی).

  \item \textbf{مداخله بدون پیوند ارگانیک} معمولاً در قالب \concept{امپریالیسم صرف} یا \concept{کشورگشایی} طبقه‌بندی می‌شود: مغول‌ها در ایران، بلژیک در کنگو، ژاپن در چین.

  \item \textbf{منطقه‌ی خاکستری:} مواردی مانند حمله‌ی امریکا به عراق (۲۰۰۳) که ترکیبی از ایدئولوژی (دموکراسی‌سازی)، منافع اقتصادی (نفت) و نجات‌طلبی بخشی از جامعه (کردها، شیعیان) بود — در هیچ‌یک از سه دسته به‌تنهایی نمی‌گنجد.
\end{itemize}
\end{tcolorbox}

% ─────────────────────────────────────────────────────────────────────────
\section{استعمارزدایی و تشکیل کشورهای جدید}
\label{sec:ch09-decolonization}
% ─────────────────────────────────────────────────────────────────────────

\subsection{تجزیه‌ی مرزها: کجا به کشور جدید انجامید؟}
\label{subsec:ch09-new-states}

فرایند استعمارزدایی (عمدتاً ۱۹۴۵–۱۹۷۵) و نیز تحولات پس از جنگ سرد، به تشکیل ده‌ها کشور جدید انجامید. پرسش کلیدی: \textbf{چه عواملی تعیین می‌کند که نجات‌طلبی/جدایی‌طلبی به تشکیل کشور مستقل بینجامد یا نه؟}

\begin{landscape}
\begin{table}[htbp]
\centering
\caption{مقایسه‌ی موارد موفق و ناموفق تشکیل کشور جدید}\label{tab:ch09-new-states}
\begin{adjustbox}{max width=\linewidth}
\begin{tabular}{
  >{\raggedleft\arraybackslash\bfseries\small}p{2.2cm}
  >{\raggedleft\arraybackslash\small}p{1.5cm}
  >{\raggedleft\arraybackslash\small}p{2cm}
  >{\raggedleft\arraybackslash\small}p{2cm}
  >{\raggedleft\arraybackslash\small}p{2.5cm}
  >{\raggedleft\arraybackslash\small}p{2cm}
  >{\raggedleft\arraybackslash\small}p{2cm}
  >{\raggedleft\arraybackslash\small}p{1.5cm}}
\toprule
\rowcolor{PrimaryDeep}
\textcolor{white}{مورد} &
\textcolor{white}{سال} &
\textcolor{white}{جدا شده از} &
\textcolor{white}{حامی خارجی} &
\textcolor{white}{عامل اصلی موفقیت/شکست} &
\textcolor{white}{شناسایی بین‌المللی} &
\textcolor{white}{ثبات کنونی} &
\textcolor{white}{نتیجه} \\
\midrule

بنگلادش &
۱۹۷۱ &
پاکستان &
هند &
نسل‌کشی + مداخله‌ی هند &
بله &
نسبی &
موفق \\
\rowcolor{NeutralLight}
اریتره &
۱۹۹۳ &
اتیوپی &
— &
۳۰ سال جنگ + رفراندوم &
بله &
استبداد &
مختلط \\
تیمور شرقی &
۲۰۰۲ &
اندونزی &
استرالیا/سازمان ملل &
رفراندوم + مداخله‌ی سازمان ملل &
بله &
شکننده &
مختلط%
\footnote{\lr{Martin, Ian. \textit{Self-Determination in East Timor}. Lynne Rienner, 2001.}} \\
\rowcolor{NeutralLight}
کوزوو &
۲۰۰۸ &
صربستان &
ناتو/امریکا &
مداخله‌ی ناتو + پاک‌سازی قومی &
ناقص &
شکننده &
مختلط \\
سودان جنوبی &
۲۰۱۱ &
سودان &
امریکا/سازمان ملل &
رفراندوم + فشار بین‌المللی &
بله &
جنگ داخلی &
شکست%
\footnote{\lr{Johnson, Douglas H. \textit{South Sudan: A New History for a New Nation}. Ohio UP, 2016.}} \\
\rowcolor{NeutralLight}
کریمه &
۲۰۱۴ &
اوکراین &
روسیه &
الحاق نظامی &
خیر (اکثریت) &
اشغال &
غیرقانونی \\
کاتالونیا &
۲۰۱۷ &
اسپانیا &
— &
رفراندوم غیرقانونی &
خیر &
— &
شکست \\
\rowcolor{NeutralLight}
کردستان عراق &
۲۰۱۷ &
عراق &
— &
رفراندوم بدون حمایت بین‌المللی &
خیر &
— &
شکست \\
چچن &
۱۹۹۴–۲۰۰۹ &
روسیه &
— (حمایت محدود) &
سرکوب نظامی &
خیر &
استبداد &
شکست%
\footnote{\lr{Lieven, Anatol. \textit{Chechnya: Tombstone of Russian Power}. Yale UP, 1998.}} \\
\rowcolor{NeutralLight}
بیافرا &
۱۹۶۷–۱۹۷۰ &
نیجریه &
فرانسه (محدود) &
شکست نظامی + قحطی &
خیر &
— &
شکست \\

\bottomrule
\end{tabular}
\end{adjustbox}
\end{table}
\end{landscape}

\subsection{عوامل تعیین‌کننده‌ی موفقیت جدایی}
\label{subsec:ch09-factors}

\begin{tcolorbox}[successbox, title={شرایطی که جدایی را ممکن می‌سازد}]
\begin{enumerate}[nosep, label=\textbf{\arabic*.}]
  \item \textbf{حمایت حداقل یک ابرقدرت:} بدون حمایت امریکا، روسیه یا اتحادیه‌ی اروپا، جدایی عملاً ناممکن است (استثنا: اریتره).
  \item \textbf{ضعف/فروپاشی دولت مادر:} بنگلادش (پاکستان ضعیف)، جمهوری‌های شوروی (فروپاشی)، یوگسلاوی (فروپاشی).
  \item \textbf{سابقه‌ی مرز اداری:} اکثر کشورهای جدید بر مبنای مرزهای اداری پیشین شکل گرفتند (\lr{uti possidetis}).%
  \footnote{\lr{Ratner, Steven R. "Drawing a Better Line: \textit{Uti Possidetis} and the Borders of New States." \textit{AJIL} 90, no.~4 (1996): 590--624.}}
  \item \textbf{فاجعه‌ی انسانی مستند:} نسل‌کشی یا پاک‌سازی قومی، همدردی بین‌المللی ایجاد می‌کند (بنگلادش، کوزوو، تیمور شرقی).
  \item \textbf{ظرفیت نهادی داخلی:} جنبش‌هایی که نهاد، رهبری و ظرفیت دولت‌سازی دارند موفق‌ترند.
\end{enumerate}
\end{tcolorbox}

\begin{tcolorbox}[failbox, title={شرایطی که جدایی را ناممکن می‌سازد}]
\begin{enumerate}[nosep, label=\textbf{\arabic*.}]
  \item \textbf{مخالفت فعال ابرقدرت‌ها:} چچن (مخالفت غرب با تضعیف روسیه)، کاتالونیا (مخالفت اتحادیه‌ی اروپا).
  \item \textbf{قدرت دولت مادر:} روسیه (چچن)، چین (تبت، اویغور)، ایران.
  \item \textbf{فقدان مرز اداری پیشین:} کردستان (تقسیم میان چهار کشور بدون مرز اداری واحد).
  \item \textbf{نگرانی از «اثر دومینو»:} اگر جدایی یک منطقه، جدایی مناطق دیگر را تشویق کند، جامعه‌ی بین‌المللی مقاومت می‌کند.
  \item \textbf{فقدان فاجعه‌ی انسانی آشکار:} بدون «بحران» واضح، حمایت بین‌المللی جلب نمی‌شود.
\end{enumerate}
\end{tcolorbox}

% ─────────────────────────────────────────────────────────────────────────
\section{اقلیت‌های تحت ستم سیستماتیک: ۵۰ سال اخیر}
\label{sec:ch09-oppressed-minorities}
% ─────────────────────────────────────────────────────────────────────────

\begin{longtable}{
  >{\raggedleft\arraybackslash\bfseries}p{2cm}
  >{\raggedleft\arraybackslash}p{2cm}
  >{\raggedleft\arraybackslash}p{2.5cm}
  >{\raggedleft\arraybackslash}p{2.5cm}
  >{\raggedleft\arraybackslash}p{2cm}
  >{\raggedleft\arraybackslash}p{2cm}}
\caption{اقلیت‌های تحت ستم سیستماتیک و واکنش بین‌المللی}\label{tab:ch09-oppressed}\\
\toprule
\rowcolor{PrimaryDeep}
\textcolor{white}{\textbf{اقلیت}} &
\textcolor{white}{\textbf{دولت ستم‌گر}} &
\textcolor{white}{\textbf{نوع ستم}} &
\textcolor{white}{\textbf{حمایت بین‌المللی}} &
\textcolor{white}{\textbf{مداخله‌ی نظامی}} &
\textcolor{white}{\textbf{نتیجه}} \\
\midrule
\endfirsthead
\toprule
\rowcolor{PrimaryDeep}
\textcolor{white}{\textbf{اقلیت}} &
\textcolor{white}{\textbf{دولت ستم‌گر}} &
\textcolor{white}{\textbf{نوع ستم}} &
\textcolor{white}{\textbf{حمایت بین‌المللی}} &
\textcolor{white}{\textbf{مداخله‌ی نظامی}} &
\textcolor{white}{\textbf{نتیجه}} \\
\midrule
\endhead
\bottomrule
\endlastfoot

اویغورها &
چین &
بازداشت جمعی، نسل‌کشی فرهنگی &
بیانیه‌ها، تحریم محدود &
خیر &
تداوم ستم%
\footnote{\lr{Zenz, Adrian. "'Thoroughly Reforming Them Towards a Healthy Heart Attitude': China's Political Re-Education Campaign in Xinjiang." \textit{Central Asian Survey} 38, no.~1 (2019): 102--128.}} \\
\rowcolor{NeutralLight}
چچن‌ها &
روسیه &
سرکوب نظامی، ناپدیدسازی &
حداقلی &
خیر &
استبداد کادیروف \\
ایزدی‌ها &
داعش &
نسل‌کشی، بردگی جنسی &
مداخله‌ی هوایی (ائتلاف) &
بله (محدود) &
بقا اما آوارگی%
\footnote{\lr{UN Human Rights Council. "They Came to Destroy." A/HRC/32/CRP.2, 2016.}} \\
\rowcolor{NeutralLight}
روهینگیا &
میانمار &
پاک‌سازی قومی، نسل‌کشی &
بیانیه‌ها &
خیر &
آوارگی گسترده \\
تیموری‌ها &
اندونزی &
اشغال، سرکوب &
رفراندوم + \lr{INTERFET} &
بله &
استقلال%
\footnote{\lr{Martin, 2001.}} \\
\rowcolor{NeutralLight}
کردها (عراق) &
بعث عراق &
انفال، نسل‌کشی &
منطقه‌ی پرواز ممنوع &
بله (غیرمستقیم) &
خودمختاری \\
دارفوری‌ها &
سودان &
نسل‌کشی &
\lr{UNAMID} &
نیمه &
تداوم بحران%
\footnote{\lr{Flint, Julie, and Alex de Waal. \textit{Darfur: A New History of a Long War}. Zed, 2008.}} \\
\rowcolor{NeutralLight}
تبتی‌ها &
چین &
سرکوب فرهنگی-مذهبی &
حمایت نمادین &
خیر &
تداوم ستم \\

\end{longtable}

\begin{tcolorbox}[enemybox, title={نتیجه‌ی تلخ: گزینشی‌بودن حمایت بین‌المللی}]
بررسی موارد بالا نشان می‌دهد که حمایت بین‌المللی از اقلیت‌های تحت ستم \textbf{کاملاً گزینشی} است:
\begin{itemize}[nosep, rightmargin=1em]
  \item \textbf{ایزدی‌ها:} مداخله شد — زیرا داعش دشمن مشترک بود و منافع ژئوپولیتیک ایجاب می‌کرد.
  \item \textbf{اویغورها:} مداخله نشد — زیرا چین ابرقدرت اقتصادی-نظامی است.
  \item \textbf{روهینگیا:} مداخله نشد — زیرا میانمار اهمیت ژئوپولیتیک کمی دارد و چین حامی آن است.
  \item \textbf{قاعده:} مداخله‌ی انسان‌دوستانه فقط زمانی رخ می‌دهد که \textbf{منافع ژئوپولیتیک قدرت‌های بزرگ} ایجاب کند — نه صرفاً وجود رنج بشری.
\end{itemize}
\end{tcolorbox}

% ─────────────────────────────────────────────────────────────────────────
\section{جمع‌بندی فصل}
\label{sec:ch09-summary}
% ─────────────────────────────────────────────────────────────────────────

\begin{tcolorbox}[wavebox, title={یافته‌های کلیدی فصل ۹}]
\begin{enumerate}[nosep, label=\textbf{\arabic*.}]
  \item استعمار بدون \concept{تسهیل‌گری محلی} ممکن نبود — اما تسهیل‌گری محلی «نجات‌طلبی» نیست بلکه «فرصت‌طلبی نخبگان» است.

  \item فعالیت‌های میسیونری، پان‌ایدئولوژیک و امپریالیستی سه شکل متفاوت «تسهیل مداخله» هستند که گاه با هم ترکیب می‌شوند.

  \item \textbf{نجات‌طلبی واقعی عمدتاً در چارچوب پیوندهای ارگانیک} (مذهب، زبان، نژاد، ایدئولوژی) رخ می‌دهد. مداخله بدون پیوند ارگانیک معمولاً امپریالیسم صرف است.

  \item تشکیل کشور جدید مستلزم ترکیبی از عوامل است: حمایت ابرقدرت، ضعف دولت مادر، مرز اداری پیشین، فاجعه‌ی انسانی و ظرفیت نهادی.

  \item حمایت بین‌المللی از اقلیت‌های تحت ستم \textbf{کاملاً گزینشی} است و تابع منافع ژئوپولیتیک — نه اصول اخلاقی.
\end{enumerate}
\end{tcolorbox}

\cleardoublepage